\documentclass[12pt]{article}
\usepackage[utf8]{inputenc}
\usepackage[T1]{fontenc}
\usepackage{hyperref}
\title{Zia Language Resources for K-12 Curriculum Development}
\author{Compiled by Claude Sonnet 5.5 (claude-sonnet-5-5) for Andr\'as Kornai}
\date{}
\begin{document}
\maketitle

\section{Introduction}
Zia (ISO 639-3 zia, Glottocode ziaa1250) is a Trans-New Guinea language of the Binandere subgroup, spoken in the lower Waria Valley, Morobe Province, Papua New Guinea. Wikipedia gives about 4,500 speakers (2000); other sources give roughly 6,000 or more. It is not known to be taught in schools. Little published material was found.

\section{Primary Educational Resources}
No primers or textbooks were found in search results.

\section{Linguistic and Reference Materials}
\begin{itemize}
\item Wikipedia article on the Zia language: \url{https://en.wikipedia.org/wiki/Zia_language}
\item Glottolog reference for Darryl Wilson (1980), an unpublished 71-page manuscript, \emph{A brief comparative grammar of Zia and Suena}: \url{https://glottolog.org/resource/reference/id/302428}
\item Glottolog reference for Darryl Wilson, \emph{The Zia orthography} (1981): \url{https://glottolog.org/resource/reference/id/313272}
\item TransNewGuinea.org Zia lexical data: \url{https://transnewguinea.org/language/zia?page=4}
\item Max Planck Institute (MPI) page on the Zia number system: \url{https://lingweb.eva.mpg.de/channumerals/Zia.htm}
\item Wikitongues language page: \url{https://wikitongues.org/languages/zia/}
\end{itemize}
No dedicated Zia dictionary was found.

\section{Religious and Community Resources}
The Zia New Testament (1982) and Old Testament work are reported. The Bible listing found in search results: \url{https://dbs.org/bibles/ZIAWBT}

\section{Digital and Audio Resources}
PARADISEC catalogue item (an archive record; its content was not examined): \url{https://catalog.paradisec.org.au/repository/AC2/FNBPNGEH109}

\section{Curriculum Development Context}
Zia has an orthography and a published New Testament, which could serve as the first readers. A curriculum would need new primers and a dictionary. The Wilson orthography work is the starting point. Community consultation is needed first.

\end{document}